\section{Score Manipulation：实用的分数操控}

第一类方法是直接操控去噪过程中的分数函数。这类方法没有严格的理论解释，但在实践中往往效果不错，尤其在图像和视频领域已有大量成功案例。

\subsection{RePaint：图像修复的经典案例}

\textbf{RePaint} 是分数操控的一个典型例子，它解决的是\textbf{图像修复（Inpainting）}问题：给定一张部分被遮挡的图像，填补缺失区域，使结果自然且与已知背景一致。

\begin{importantbox}{RePaint 的核心思想}
在每个去噪步骤中，将中间数据点 $x_t$ 分为两部分处理：
\begin{itemize}
    \item \textbf{前景区域}（待填补区域）：执行正常的\textbf{逆过程}（reverse process），从噪声中逐步生成内容。
    \item \textbf{背景区域}（已知区域）：执行\textbf{前向过程}（forward process），从已知的干净图像加噪到对应时间步 $t-1$ 的噪声水平。
\end{itemize}
然后将两部分拼合，得到 $x_{t-1}$，重复此过程直到 $t=0$。
\end{importantbox}

具体算法流程如下：
\begin{enumerate}
    \item 给定当前中间数据 $x_t$，执行一步逆过程，得到前景区域的 $x_{t-1}^{\text{fg}}$。
    \item 对已知的干净背景图像执行一步前向过程到时间步 $t-1$，得到 $x_{t-1}^{\text{bg}}$。
    \item 用掩码（mask）将两者组合：前景取 $x_{t-1}^{\text{fg}}$，背景取 $x_{t-1}^{\text{bg}}$。
    \item 重复上述过程。
\end{enumerate}

\begin{warningbox}{RePaint 缺乏理论保证}
为什么背景能被保留是直观的——我们在每一步都从干净图像重新加噪。但前景和背景之间的\textbf{无缝衔接}并没有理论保证。实际使用中确实存在前景与背景不匹配的失败案例。讲者强调：``No theory here, just practical ideas.''
\end{warningbox}

\subsection{Attention Map 操控}

除了直接操控分数/噪声外，还可以通过操控噪声预测网络内部的\textbf{注意力图}（Attention Maps）来实现引导，包括：

\begin{itemize}
    \item \textbf{Self-Attention 操控}：修改自注意力模块中的注意力权重，影响图像内部不同区域之间的关联。
    \item \textbf{Cross-Attention 操控}：修改交叉注意力中文本 token 与图像区域的对应关系，从而控制特定对象的位置和属性。
\end{itemize}

\begin{knowledgebox}{Cross-Attention 操控的应用示例}
通过操控交叉注意力，可以实现以下任务（均无需重新训练扩散模型）：
\begin{itemize}
    \item \textbf{风格迁移}：从一张图像提取结构信息，从另一张提取风格/外观信息，生成新图像。
    \item \textbf{布局控制}：给定文本描述和 2D 边界框，控制各对象在图像中的位置。例如，指定``黄色汽车''在右侧黄色框中，``城堡''在左侧绿色框中。
    \item \textbf{视频编辑}：通过反转（inversion）和注意力操控，将视频中的对象移动到新位置，同时保持视频的自然性。
\end{itemize}
\end{knowledgebox}

\subsection{Score Manipulation 的优缺点}

\begin{table}[H]
\centering
\begin{tabular}{p{0.45\textwidth}p{0.45\textwidth}}
\toprule
\textbf{优点} & \textbf{缺点} \\
\midrule
直接、直观、易于实现 & 没有理论保证 \\
在图像/视频领域有大量成功案例 & 存在失败案例，结果不可预测 \\
无需额外训练或数据 & 多为针对特定任务设计的 ad-hoc 方法 \\
可拓展到其他数据模态 & 难以泛化到新的引导类型 \\
\bottomrule
\end{tabular}
\caption{Score Manipulation 方法的优缺点对比}
\end{table}

\subsection{本章小结}

Score Manipulation 是一类高度实用的推理时引导技术。它的核心思想是在去噪的每一步中，直接修改分数函数或注意力图，使生成轨迹偏向期望的方向。虽然缺乏理论保证，但在图像修复、风格迁移、布局控制、视频编辑等任务上已有大量成功案例。接下来我们将讨论一类更加系统化、具有理论基础的方法——基于逆问题的引导。


